% from ega0-21-fr.tex
\oldpage[0\textsubscript{IV-1}]{171}
Par suite, on a l'inégalité
\begin{equation}
\tag{21.7.3.2}
\label{0.21.7.3.2-fr}
\operatorname{rg}_L\Omega^1_{L/k}-\operatorname{rg}_K\Omega^1_{K/k}
\geqslant\deg.\operatorname{tr}_K L.
\end{equation}
L'égalité étant atteinte si et seulement si $k$ est un corps admissible pour
l'extension $L$ de $K$.
\end{corollary}

\begin{proof}
On a en effet la suite exacte (20.6.1.1)
\[
0\to\Upsilon_{L/K/k}\to\Omega^1_{K/k}\otimes_KL\to\Omega^1_{L/k}
\to\Omega^1_{L/K}\to0
\]
donc le premier membre de (21.7.3.1) est égal à
\[
\operatorname{rg}_L\Omega^1_{L/k}-\operatorname{rg}_L\Upsilon_{L/K/k}
\]
et il suffit d'appliquer (21.7.2).
\end{proof}

L'intérêt de ce dernier corollaire est qu'il ne fait intervenir que des
modules de différentielles, à l'exclusion de modules d'imperfection. Nous
verrons d'ailleurs plus bas (21.8.6) que pour toute extension de type fini
$L$ de $K$, il existe un sous-corps $k$ de $K$ tel que $[K:k]$ soit fini et
qui est admissible pour l'extension $L$ de $K$, de sorte que les deux membres
de (21.7.3.2) sont alors égaux.

\begin{corollary}[21.7.4]
\label{0.21.7.4-fr}
Soit $K$ une extension de type fini d'un corps $k$. Les conditions suivantes
sont équivalentes~:
\begin{enumerate}
  \item[\rm a)] $K$ est une extension finie séparable de $k$.
  \item[\rm b)] $K$ est une $k$-algèbre formellement non ramifiée (19.10.2).
  \item[\rm c)] $K$ est une $k$-algèbre formellement étale (19.10.2).
  \item[\rm d)] On a $\Omega^1_{K/k}=0$.
\end{enumerate}
\end{corollary}

\begin{proof}
En effet, on sait (les topologies étant discrètes) que les conditions b) et
d) sont équivalentes (20.7.4) et que a) entraîne que $K$ est une
$k$-algèbre formellement lisse (19.6.1), donc la conjonction de a) et b)
équivaut à c). Tout revient à voir que d) entraîne a). Or, en vertu de
(21.7.1.1), la relation $\Omega^1_{K/k}=0$ entraîne que $K$ est algébrique
(donc fini) sur $k$ et que $\Upsilon_{K/k}=0$, c'est-à-dire (20.6.3) que
$K$ est séparable sur $k$.
\end{proof}

\begin{remark}[21.7.5]
\label{0.21.7.5-fr}
En vertu de (20.6.3.2) et de (0$_{\rm III}$, 11.10.2), le premier membre de
(21.7.1.1) n'est autre que la \emph{caractéristique d'Euler-Poincaré} du
complexe $K_\bullet(C/B/A)$ introduit dans (20.6.3), pour $C=L$, $B=K$ et
$A=P$ (corps premier de $K$). Dans le chapitre consacré à la théorie des
intersections et au théorème de Riemann--Roch, un rôle important sera joué
également par les caractéristiques d'Euler-Poincaré généralisées (à valeurs
dans des groupes de classes de $\mathcal O_X$-Modules) de complexes
généralisant les complexes $F_\bullet(C/A)$ considérés dans (20.6.22).
\end{remark}

\subsection{Critères d'admissibilité}
\label{subsection:0.21.8-fr}

Nous revenons à nos conventions antérieures et supposons donc que tous les
corps considérés dans ce numéro sont de caractéristique $p>0$.

\begin{lemma}[21.8.1]
\label{0.21.8.1-fr}
Soient $K$ un corps, $k$ un sous-corps de $K$, $(k_\alpha)_{\alpha\in I}$ une
famille filtrante décroissante de sous-corps de $K$ telle que
$k=\bigcap_{\alpha\in I}k_\alpha$. Soient $V$ un espace vectoriel sur $K$,
$(a_i)_{1\leqslant i\leqslant n}$ une famille finie de vecteurs de $V$~; si la
famille $(a_i)$ est libre sur $k$, il existe un indice $\gamma$ tel qu'elle
soit aussi libre sur $k_\gamma$.
\end{lemma}

